% =============================================================================
% 00_opening.tex — Slides S01-001 to S01-004
% Session Opening, The Dynamic Typing Paradox, North-Star Question & Roadmap
% =============================================================================

% -----------------------------------------------------------------------------
% Slide ID: S01-001
% Narrative Role: Title / Cover
% Cognitive Objective: Establish session identity, module context, and program standards.
% Plan Source: slide_plan.md / S01-001
% -----------------------------------------------------------------------------
\UpGradTitleSlide{Python Programming Foundations}{Mastering the Python Virtual Machine, Memory Semantics \& Vectorized Mechanics}{M00: Foundations of Data Science}{Session 01}

% -----------------------------------------------------------------------------
% Slide ID: S01-002
% Narrative Role: Hook / Contrast
% Cognitive Objective: Confront the learner with the fundamental contradiction between compiled typed languages and Python.
% Plan Source: slide_plan.md / S01-002
% -----------------------------------------------------------------------------
\begin{frame}[fragile]{The Dynamic Typing Paradox}
  \begin{columns}[T,onlytextwidth]
    \begin{column}{0.48\textwidth}
      {\small\bfseries\color{UpGradRed}C++ / Java (Statically Typed Boxes)}\par\vspace{0.08cm}
      \begin{HighlightedCodeBlock}{C++}{2}
int x = 42;
x = "hello"; // FATAL COMPILATION ERROR!
      \end{HighlightedCodeBlock}
      \vspace{0.08cm}
      \begin{tcolorbox}[enhanced,arc=1mm,boxrule=.4pt,colback=UpGradRed!10,colframe=UpGradRed,left=2mm,right=2mm,top=1.5mm,bottom=1.5mm]
        {\scriptsize\bfseries\color{UpGradRedDark}Compile Error:}\enspace{\tiny\ttfamily cannot convert 'const char*' to 'int'}
      \end{tcolorbox}
      \vspace{0.04cm}
      {\scriptsize\color{UpGradSlate}Variables are fixed, statically typed physical memory slots.}
    \end{column}
    \hfill{\color{UpGradRule}\vrule width .5pt}\hfill
    \begin{column}{0.48\textwidth}
      {\small\bfseries\color{AWSBlue}Python (Dynamic Pointer Rebinding)}\par\vspace{0.08cm}
      \begin{HighlightedCodeBlock}{Python}{2,3}
x = 42
x = "hello"     # Executes seamlessly!
x = [1, 2, 3]   # No compiler complaints!
      \end{HighlightedCodeBlock}
      \vspace{0.08cm}
      \begin{tcolorbox}[enhanced,arc=1mm,boxrule=.4pt,colback=AWSEmerald!10,colframe=AWSEmerald,left=2mm,right=2mm,top=1.5mm,bottom=1.5mm]
        {\scriptsize\bfseries\color{AWSEmerald}Runtime Status:}\enspace{\tiny\bfseries\color{AWSEmerald}EXECUTES SEAMLESSLY (Exit Code 0)}
      \end{tcolorbox}
      \vspace{0.04cm}
      {\scriptsize\color{UpGradSlate}Why does Python permit one variable to effortlessly change types?}
    \end{column}
  \end{columns}
  \vspace{0.18cm}
  \TakeawayBanner{In Python, variables are not data boxes—they are named memory pointers.}
\end{frame}

% -----------------------------------------------------------------------------
% Slide ID: S01-003
% Narrative Role: North-Star Hook
% Cognitive Objective: Frame the overarching question that unifies memory semantics, closures, and NumPy performance.
% Plan Source: slide_plan.md / S01-003
% -----------------------------------------------------------------------------
\begin{frame}{The Session North-Star Question}
  \begin{center}
    \vspace{0.05cm}
    \begin{tcolorbox}[enhanced,arc=2mm,boxrule=.6pt,colback=UpGradMist,colframe=UpGradRule,
      borderline west={3.5pt}{0pt}{UpGradRed},left=4.5mm,right=4.5mm,top=3.5mm,bottom=3.5mm,width=0.96\linewidth]
      {\UpGradDisplayFont\fontsize{12}{16}\selectfont\bfseries\color{UpGradNight}
        “How can Python let the same variable appear to become an integer, a collection, or even a function wrapper, and what does that dynamic flexibility cost us when scaling numerical computations in pure Python vs NumPy?”
      }
    \end{tcolorbox}
  \end{center}
  \vspace{0.15cm}
  \begin{columns}[c,onlytextwidth]
    \begin{column}{0.31\textwidth}
      \StatCard{UpGradRed}{Block 1}{Pointer Semantics}{Stack names \& heap PyObjects}
    \end{column}
    \begin{column}{0.31\textwidth}
      \StatCard{AWSBlue}{Block 2}{Hash Tables \& Streams}{O(1) lookups \& lazy generators}
    \end{column}
    \begin{column}{0.31\textwidth}
      \StatCard{AWSEmerald}{Block 3}{Vectorized SIMD}{C-arrays, strides \& broadcasting}
    \end{column}
  \end{columns}
\end{frame}

% -----------------------------------------------------------------------------
% Slide ID: S01-004
% Narrative Role: Roadmap / Transition
% Cognitive Objective: Provide clear structural anchors for the 180-minute journey.
% Plan Source: slide_plan.md / S01-004
% -----------------------------------------------------------------------------
\SessionAgendaSlide{
  \BlockEntry{01}{The Object Pointer Mental Model \& Memory Semantics}{Stack pointers vs heap allocations, PyObject anatomy, in-place mutation, default argument leakage, and shallow vs deep copies.}
  \BlockEntry{02}{Data Structures, Lazy Generators \& Practical Decorators}{Hash table direct indexing, why mutable keys fail, O(1) memory streams with yield, first-class functions, and robust decorators with @wraps.}
  \BlockEntry{03}{The CPython Speed Wall \& NumPy's Breakthrough}{3 CPU bottlenecks of pure Python loops, contiguous C-memory layout, SIMD vectorization, strides math, and zero-memory broadcasting.}
}
